\documentclass[10pt,a4paper]{amsart}
\usepackage[margin=19mm]{geometry}
\usepackage{amsmath,amssymb,mathtools}
\usepackage[scr=rsfs,cal=euler]{mathalfa}
\usepackage{xfrac}
\usepackage[unicode,hidelinks,bookmarks=false]{hyperref}
\newcommand{\ie}{i.e.\ }
\newcommand{\cf}{cf.\ }
\newcommand{\resp}{respectively\ }
\newcommand{\Subsection}{\subsection}
\providecommand{\nobreakdash}{\nobreak\hbox{-}}
\setlength{\emergencystretch}{2em}
\multlinegap0pt
\usepackage{bm}
\usepackage{bbm}
\usepackage{dsfont}
\usepackage{xspace}
\usepackage{cancel}
\usepackage{mathtools}
\usepackage{amsthm}
\makeatletter
\@ifundefined{theorem}{\newtheorem{theorem}{Theorem}}{}
\@ifundefined{lemma}{\newtheorem{lemma}[theorem]{Lemma}}{}
\makeatother
\DeclareMathOperator*{\argmax}{arg\,max}
\title[On Theoretical Interpretations of Concept-Based In-Context L]{On Theoretical Interpretations of Concept-Based In-Context Learning}
\author{Huaze Tang, Tianren Peng, Shao-lun Huang}
\date{Source: arXiv:2509.20882v2 (CC BY 4.0)}
\begin{document}
\maketitle

\noindent\emph{Source and licence.} On Theoretical Interpretations of Concept-Based In-Context Learning. Version arXiv:2509.20882v2 (CC BY 4.0), distributed under the Creative Commons Attribution 4.0 International licence, as recorded in the arXiv licence metadata for this exact version and shown on the abstract page https://arxiv.org/abs/2509.20882v2. Full rights notice: licenses/arxiv-2509.20882-CC-BY-4.0.txt.\par\smallskip

\noindent\textit{Editorial scope.} Counted unit: the Lemma whose source label is \texttt{lem:upper-bound} in the article (source0.tex lines 991--996, source label \texttt{lem:upper-bound}), delivered together with the article's own complete original proof of it (source0.tex lines 997--1043, one \texttt{proof} environment of 47 lines). No mathematics is written, repaired or re-derived: statement and proof are the article's own text line for line. Required context retained uncounted: equ:eig-root (lines 968--970); equ:spectrum\_equ (lines 975--977). Every definition, display and label of those lines is retained unchanged. Omitted from this package and not redistributed: 1--967, 971--974, 978--990, 1044--1833. Nothing inside a retained excerpt is omitted or altered: every retained body is delivered exactly as the article's lines are written, so no omission receipt of any kind accompanies this package. Citations: the retained excerpts carry no citation command at all, so no bibliography entry of the article is transcribed and no reference is redistributed. Reference adaptation: none was needed and none was made; every reference inside the delivered unit resolves inside it, so no unresolved reference or citation is produced. Printed slips kept literal: no printed word, symbol, hypothesis or number of the counted statement or of its proof was altered. Layout and class: the article's own class is \texttt{article} with packages including iclr2026\_conference, times, amsmath, amsfonts, bm, hyperref, graphicx, url, amssymb, amsthm, booktabs, multirow, siunitx, makecell; the wrapper is the lane's pinned \texttt{amsart} editorial wrapper at 10pt on A4, which restates the theorem-like environments the retained text needs and adds the article's own macro definitions that the retained text needs (\texttt{\textbackslash argmax}), with extra wrapper packages (bm, bbm, dsfont, xspace, cancel, mathtools). The delivered numbering is the unit's own. Assets: no retained span contains a figure environment, a graphics-inclusion command, a PDF-page inclusion, a table or a TikZ picture; no image file is copied into this package, the package declares no asset dependency and no image file is redistributed with it. Licence: version arXiv:2509.20882v2 of this article is distributed under the Creative Commons Attribution 4.0 International licence, as recorded in the arXiv licence metadata for this exact version and shown on the abstract page https://arxiv.org/abs/2509.20882v2. A full rights notice is delivered at licenses/arxiv-2509.20882-CC-BY-4.0.txt. Characters: every retained excerpt is pure ASCII, so the delivered engine, XeLaTeX with its default Latin Modern text font, reports no missing character.
\begin{equation}\label{equ:eig-root}
1+\sum_{j=1}^M \frac{P_{Y|X}^2(j|x_i)}{\lambda - P_{Y|X}(j|x_i)}=0,
\end{equation}

\begin{equation}\label{equ:spectrum_equ}
\det(\lambda I- \mathsf{diag}\left\{ P_{Y|X=x_i} \right\} + \phi_i\phi_i^\mathrm{T})=0, 
\end{equation}

\begin{lemma}\label{lem:upper-bound}
For each $i$, the largest eigenvalue of $\textbf{A}(x_i)$ satisfies
\[
\lambda_1(\textbf{A}(x_i)) \leq 2 \max_j P_{Y|X}(j|x_i)\big(1-\max_j P_{Y|X}(j|x_i)\big).
\]
\end{lemma}

\begin{proof}
To prove this result, we distinguish between two cases: whether the maximum probability is attained by more than one label, or by a unique label.

\paragraph{Case 1: The size of set $\arg\max_j P_{Y|X}(j|x_i)$ is larger than $1$}
By continuity, we have 
\begin{equation*}
\det(P_{Y|X}(j|x_i)\cdot I- {\rm diag}\left\{ P_{Y|X=x_i} \right\} + \phi_i\phi_i^\mathrm{T})= P_{Y|X}^2(j|x_i)\prod_{k\ne j}(P_{Y|X}(j|x_i)-P_{Y|X}(k|x_i)).
\end{equation*}
Since $\arg\max_y P_{Y|X}(y|x_i)$ is not a single point set, $\max_y P_{Y|X}(y|x_i)$ is one solution to Eq. (\ref{equ:spectrum_equ}). Then, the largest eigenvalue of $\mathbf{A}(x_i)$ is $\max_y P_{Y|X}(y|x_i)$. In this case, since the size of set  $\arg\max_j P_{Y|X}(j|x_i)$ is larger than $1$, we have that $\max_j P_{Y|X}(j|x_i) \leq 1/2$, which gives that $\max_j P_{Y|X}(j|x_i) \leq 2\max_j P_{Y|X}(j|x_i)(1-\max_j P_{Y|X}(j|x_i))$.

\paragraph{Case 2: The size of set  $\arg\max_j P_{Y|X}(j|x_i)$ equals to $1$}
Let $j_0=\arg\max_j P_{Y|X}(j|x_i)$ and denote the largest root of the Eq. (\ref{equ:eig-root}) by $\lambda_0$.
Further, let arbitrary $j_1\in\argmax_{j,j\ne j_0} P_{Y|X}(j|x_i)$. Note that
\begin{equation*}
    1+\sum_{j=1}^M\frac{P_{Y|X}^2(j|x_i)}{\lambda-P_{Y|X}(j|x_i)}>1,\forall\lambda>P_{Y|X}(j_0|x_i),
\end{equation*}
\begin{equation*}
    \lim_{\lambda\to P_{Y|X}(j_0|x_i)^-}1+\sum_{j=1}^M\frac{P_{Y|X}^2(j|x_i)}{\lambda-P_{Y|X}(j|x_i)}=-\infty,
\end{equation*}
and
\begin{equation*}
    \lim_{\lambda\to P_{Y|X}(j_1|x_i)^+}1+\sum_{j=1}^M\frac{P_{Y|X}^2(j|x_i)}{\lambda-P_{Y|X}(j|x_i)}=\infty.
\end{equation*}
We have $P_{Y|X}(j_1|x_i) < \lambda_0 <P_{Y|X}(j_0|x_i)$. Following the analysis in Case 1, the value $P_{Y|X}(j_0|x_i)$ will not be the root of Eq. (\ref{equ:spectrum_equ}). Hence, we have the largest eigenvalue of $\mathbf{A}(x_i)$ must be $\lambda_0$.

In this case, we have the
\begin{equation*}
    \frac{P_{Y|X}^2(j_0|x_i)}{\lambda_0-P_{Y|X}(j_0|x_i)}+\frac{P_{Y|X}^2(j_1|x_i)}{\lambda_0-P_{Y|X}(j_1|x_i)}\le\overbrace{\sum_{j=1}^M\frac{P_{Y|X}^2(j|x_i)}{\lambda_0-P_{Y|X}(j|x_i)}}^{=-1}
\end{equation*}
and
\begin{align*}
    \overbrace{\sum_{j=1}^M\frac{P_{Y|X}^2(j|x_i)}{\lambda_0-P_{Y|X}(j|x_i)}}^{=-1} & \le\frac{P_{Y|X}^2(j_0|x_i)}{\lambda_0-P_{Y|X}(j_0|x_i)}+\frac{\sum_{k\ne j}P_{Y|X}^2(k|x_i)}{\lambda_0-P_{Y|X}(j_1|x_i)}\\
    & \le\frac{P_{Y|X}^2(j_0|x_i)}{\lambda_0-P_{Y|X}(j_0|x_i)}+\frac{(1-P_{Y|X}(j_0|x_i))^2}{\lambda_0-(1-P_{Y|X}(j_0|x_i))}
\end{align*}

Hence, solving 
$$
\frac{P_{Y|X}^2(j_0|x_i)}{\lambda_0-P_{Y|X}(j_0|x_i)}+\frac{(1-P_{Y|X}(j_0|x_i))^2}{\lambda_0-(1-P_{Y|X}(j_0|x_i))}\ge-1
$$
gives $\lambda_0\le2P_{Y|X}(j_0|x_i)(1-P_{Y|X}(j_0|x_i))$, where the equality is achieved by Bernoulli distribution.

In conclusion of the two cases, both cases give that
\begin{equation*}
    \lambda_1(\textbf{A}(x_i)) \leq 2\max_y P_{Y|X}(y|x_i)\left(1-\max_y P_{Y|X}(y|x_i)\right).
\end{equation*}
This completes the proof.
\end{proof}

\end{document}
